%% kapa: exam=1 round=30 year=2019 period=1 form=A subject=real_estate subject_ko=부동산학원론 seq=12 number=92
%% source: Q-Net 감정평가사 기출문제 (cst00309) artlSeq=5205205 file=extracted/1008/5205205/2019년 제30회 감정평가사 1차시험 문제지/2019년 제30회 감정평가사 1차시험 문제지/2019년도 제30회 감정평가사 1차 1교시 A형.hwp
%% conversion: kordoc 4.24.51 -> tools/gen1.py (auto); review=proofread-sonnet; flags=none
\begin{question}{92}
부동산투자분석에 관한 설명으로 옳지 않은 것은?
\begin{choices}
  \item 순현재가치는 장래 예상되는 현금유입액과 현금유출액의 현재가치를 차감한 금액이다.
  \item 내부수익률은 장래 예상되는 현금유입액과 현금유출액의 현재가치를 같게 하는 할인율이다.
  \item 회수기간법은 투자안 중에서 회수기간이 가장 단기인 투자안을 선택하는 방법이다.
  \item 순현가법, 내부수익률법, 수익성지수법은 현금흐름을 할인하여 투자분석을 하는 방법이다.
  \item 순현재가치가 1보다 큰 경우나 내부수익률이 요구수익률보다 큰 경우에는 투자하지 않는다.
\end{choices}
\end{question}
\answer{5}
